% =====================================================================
%  Introduction + Related Work
%  Ultra-High-Resolution Defocus Deblurring via Blur-Guided
%  Generative Upsampling
% ---------------------------------------------------------------------
%  Macros (move to preamble if preferred)
% =====================================================================
\providecommand{\method}{\textsc{Method}}          % TODO: final name
\providecommand{\todo}[1]{{\color{red}[TODO: #1]}}
\providecommand{\etal}{\textit{et al.}}
\providecommand{\ie}{\textit{i.e.}}
\providecommand{\eg}{\textit{e.g.}}
\providecommand{\vs}{\textit{vs.}}
% Requires: \usepackage{xcolor}

% =====================================================================
\section{Introduction}
\label{sec:intro}
% =====================================================================

Defocus blur is an unavoidable consequence of finite depth of field.
Whenever a scene spans a range of depths larger than the depth of field of
the lens, regions away from the focal plane are imaged as a spatially varying
superposition of circles of confusion (CoC) whose size grows with the
distance from the focal plane and with the aperture diameter.
Wide apertures are routinely used to gather light in low-light conditions,
to enable fast shutter speeds, or simply because the photographer mis-focused,
and the result is a photograph in which only a thin slab of the scene is sharp.
Single-image defocus deblurring aims to recover an all-in-focus image from
such an observation, and has progressed rapidly with learning-based
methods~\cite{abuolaim2020dpdd,son2021kpac,lee2021ifan,ruan2022drbnet,
quan2023nrknet,zamir2022restormer,li2023jdrl}.

At the same time, the resolution of the images people actually capture has
grown far beyond the resolution at which these methods are developed and
evaluated. Interchangeable-lens cameras routinely record $24$--$60$\,MP, and
smartphone main sensors reach $50$--$200$\,MP.
In contrast, the de facto protocol in defocus deblurring trains networks on
$512\times512$ crops and evaluates on images downsampled to roughly
$2$\,MP; the standard DPDD benchmark~\cite{abuolaim2020dpdd}, for instance,
is distributed and evaluated at $1680\times1120$, although it was captured
at the full $6720\times4480$ resolution of a Canon EOS~5D Mark~IV.
% VERIFY: DPDD capture resolution / camera / released resolution.
This mismatch is not a cosmetic one. Downsampling by a factor $s$ shrinks
every CoC by the same factor, so a model evaluated at low resolution solves
a substantially easier problem: blur kernels are small relative to its
receptive field and to the training crop, and the high-frequency content
that is hardest to restore has already been discarded.
The question of how to produce an all-in-focus image \emph{at the native
resolution of the sensor} has, to our knowledge, not been addressed.

\begin{figure}[t]
  \centering
  % \includegraphics[width=\linewidth]{figures/teaser.pdf}
  \fbox{\parbox[c][4cm][c]{0.95\linewidth}{\centering Teaser placeholder}}
  \caption{\textbf{Defocus deblurring at native resolution.}
  On a $30$\,MP DPDD image, applying a state-of-the-art deblurring network
  patch-wise at native resolution leaves large-CoC regions unresolved and
  produces visible seams (a); deblurring at low resolution and upsampling
  with a super-resolution model yields inconsistent, hallucinated textures
  across the image (b). \method{} upsamples the low-resolution deblurred
  estimate under the guidance of the native-resolution blurry input,
  transferring real detail where it survives and synthesizing it where it
  does not, with related regions across the image restored consistently
  through low-resolution-routed cross-patch attention (c).
  \todo{final figure and panel labels}}
  \label{fig:teaser}
\end{figure}

Two straightforward strategies exist for applying current methods to
ultra-high-resolution (UHD) inputs, and both fail
(Fig.~\ref{fig:teaser}).
The first is to run the network directly at native resolution, necessarily
in overlapping patches because of memory constraints.
This suffers from two compounding problems.
First, the CoC of strongly defocused regions grows proportionally with
resolution and can exceed the patch size itself, so each patch lacks the
spatial context required to invert its blur; networks trained at low
resolution have furthermore never observed blur of this extent and degrade
sharply out of distribution.
Second, independent per-patch predictions disagree near patch boundaries,
producing seams, tone shifts, and texture discontinuities that blending in
overlapping regions only partially hides~\cite{chu2022tlc}.
The second strategy is to downsample the input so that the blur falls back
into the training distribution, deblur at low resolution, and upsample the
result, either with an interpolation filter or with a (generative)
super-resolution model~\cite{liang2021swinir,wang2021realesrgan,
wang2024stablesr,yu2024supir,wu2024osediff}.
Interpolation produces a soft image that no longer has native-resolution
detail. Super-resolution models reintroduce high frequencies, but they do so
by hallucination: they are blind to the native-resolution observation, and,
since they are themselves applied in tiles at these resolutions, they
synthesize textures that are inconsistent across the image and often
unrelated to the actual scene.

Our key observation is that the native-resolution blurry input is far from
uninformative, and that the difficulty of defocus deblurring is naturally
split across scales.
Because defocus blur is spatially varying, the native-resolution observation
contains, at every pixel, a different amount of recoverable information:
in-focus regions retain their full detail; mildly defocused regions still
carry frequencies above the Nyquist limit of the downsampled image; and only
strongly defocused regions have lost high-frequency content altogether
and require genuine synthesis.
Conversely, the component of the problem that is hard at native
resolution---inverting blur whose extent exceeds any tractable patch---is
exactly the component that is easy at low resolution, where the whole image
fits in a single forward pass and existing methods already perform well.
This suggests decoupling the problem: let a low-resolution deblurring
method establish the globally consistent structure of the all-in-focus
image, and let a native-resolution model add high-frequency detail, copying
it from the observation wherever it survives and generating it only where
it does not.

Building on this observation, we introduce \method{}, a \emph{blur-guided
generative upsampler} for UHD defocus deblurring.
\method{} consists of two stages connected by a compact interface.
\emph{Stage~1} operates at $4\times$ lower resolution, where the entire
image is processed jointly, and provides three signals:
(i)~a low-resolution all-in-focus \emph{anchor};
(ii)~a \emph{deblur-level map} that jointly encodes, per region, how much
restoration was required and how reliable the resulting anchor is; this
identifies where the native-resolution observation has lost detail and
generative synthesis will be needed, and where the anchor should be
trusted; and
(iii)~a whole-image \emph{affinity map} derived from self-attention,
indicating which regions of the image depict related content---nearby or
distant repetitions of the same material, surface, or structure.
When Stage~1 is our own transformer-based deblurring network, (ii) and
(iii) are read out of the model implicitly, from its internal activations
and attention maps, at negligible extra cost.
\emph{Stage~2} is a native-resolution generative upsampler that operates
patch-wise and consumes the three signals together with the
native-resolution defocused image.
The affinity map routes \emph{cross-patch attention}: each patch attends
not only to its spatial neighbors but to the patches Stage~1 identifies as
related, and overlapping predictions are fused in feature space, so that
related content is restored consistently across the entire image even
though inference runs patch-wise with bounded memory.
The deblur-level map drives a \emph{blur-aware copy/generate} mechanism
that transfers observed high-frequency detail where it survives and
synthesizes it from the generative prior only where it does not. This
grounds the output in real image content wherever possible, which is
precisely what distinguishes our setting from super-resolution.
The upsampler is trained as a multi-step conditional diffusion model and
then distilled into a one-step generator, which makes inference on tens of
megapixels practical and which, as we discuss in Sec.~\ref{sec:related},
requires cross-patch consistency to be established within a single forward
pass.

A central property of this design is that Stage~2 depends on Stage~1 only
through the interface, not through its internals. This makes \method{} a
plug-in module that lifts \emph{any} existing low-resolution deblurring
method to native resolution without retraining it. Off-the-shelf methods
provide only the anchor, and we supply the remaining two signals with
method-agnostic surrogates.
The deblur-level map is approximated by a training-free
\emph{re-blur agreement} surrogate: we re-blur the anchor with a sweep of
defocus kernels and, per region, find the kernel that best reproduces the
downsampled input. The size of that kernel estimates how much deblurring
was performed, and the residual error at that kernel measures how
consistent, and hence how reliable, the anchor is with the observation.
Combined with the local sharpness of the input, which exposes regions the
deblurring method left unrestored, these measurements are mapped to the
format of the native deblur-level map by a tiny calibration head that is
trained once, on anchors from a pool of methods, and itself depends only on
black-box quantities.
As a physics-only alternative, we also derive the map from monocular
depth~\cite{yang2024depthanything,ke2024marigold}, since the CoC grows with
the distance to the focal plane, which we locate from the sharpest regions
of the input.
Surrogates are necessarily cruder than the learned map; we quantify this
gap explicitly and show that \method{} remains effective across methods
nonetheless.
The affinity map is approximated by cross-attention between features of a
frozen vision foundation model~\cite{oquab2024dinov2,kirillov2023sam},
whose representations are known to encode reliable semantic and texture
correspondences~\cite{amir2021deepvit,tang2023dift}.
To make Stage~2 robust to which variant of each signal it receives, we
randomly interchange native and surrogate signals during training.
For the same reason, the upsampler must not overfit to the artifact
signature of any particular deblurrer. We therefore train it on a mixture
of anchors: outputs of a pool of existing deblurring methods, and
\emph{simulated} anchors obtained by degrading the downsampled ground truth
with the error modes typical of low-resolution deblurring---residual
defocus correlated with the defocus map, ringing, over-sharpening, and tone
shifts.
Finally, since real aperture-bracketed pairs at native resolution cover
only a few hundred scenes, we complement them with defocused images
synthesized from sharp high-resolution photographs using estimated depth
and a physically based CoC rendering model, and fine-tune on real pairs to
close the synthetic-to-real gap.

Evaluating at native resolution raises its own difficulties, which we
address with a dedicated protocol on the full-resolution DPDD
data~\cite{abuolaim2020dpdd}.
The all-in-focus ground truth is captured with a small aperture and is
therefore diffraction-limited at native resolution, and the blurry and
sharp captures, taken in succession, exhibit residual misalignment that is
amplified by the resolution increase~\cite{li2023jdrl}.
We therefore register image pairs before evaluation, complement
full-reference fidelity metrics with perceptual and no-reference measures,
and introduce metrics that directly quantify discontinuities along patch
boundaries and texture consistency across patches, which whole-image PSNR
largely fails to capture.
\todo{Add additional native-resolution test set(s) once finalized,
\eg, high-resolution smartphone captures.}

In summary, our contributions are:
\begin{itemize}
  \item We identify and analyze the gap between the resolution at which
        defocus deblurring is studied and the resolution at which images are
        captured, showing that both native-resolution patch-wise inference
        and low-resolution deblurring followed by super-resolution fail, for
        complementary reasons rooted in the scaling of the circle of
        confusion with resolution.
  \item We propose \method{}, a two-stage framework whose low-resolution
        stage provides an all-in-focus anchor, a deblur-level map, and a
        whole-image affinity map, and whose native-resolution stage is a
        blur-guided generative upsampler distilled to a single step. The
        affinity map routes cross-patch attention to related content
        anywhere in the image, and the deblur-level map drives a blur-aware
        copy/generate mechanism that transfers real detail where it survives
        and synthesizes it where it does not.
  \item We make the upsampler agnostic to the low-resolution method through
        method-agnostic surrogates for the interface signals---a physically
        grounded re-blur agreement measure that jointly estimates deblur
        level and anchor reliability, and foundation-feature affinities---
        together with training on a mixture of real and simulated anchors.
        \method{} consistently improves \todo{N} existing methods when
        lifting them to $30$\,MP, without retraining them, including methods
        held out from training entirely.
  \item We establish a native-resolution evaluation protocol for defocus
        deblurring on DPDD, including pair registration and patch-seam and
        cross-patch consistency metrics. \method{} outperforms native
        patch-wise inference and deblur-then-super-resolve pipelines by
        \todo{X}\,dB PSNR and \todo{Y} LPIPS while running in \todo{Z}\,s
        per $30$\,MP image.
\end{itemize}


% =====================================================================
\section{Related Work}
\label{sec:related}
% =====================================================================

\paragraph{Defocus deblurring.}
Classical approaches treat defocus deblurring as a two-step problem: a
spatially varying defocus map is first estimated, from edge
cues~\cite{zhuo2011defocusmap,karaali2018edgedefocus} or learned
predictors~\cite{lee2019dmenet}, and the image is then restored by
non-blind deconvolution~\cite{krishnan2009hyperlaplacian}.
Errors in the defocus map and the limited expressiveness of parametric
kernels make these pipelines fragile on real photographs.
Abuolaim and Brown~\cite{abuolaim2020dpdd} introduced DPDD, the first
dataset of real defocused/all-in-focus pairs captured with a large and a
small aperture, together with an end-to-end network, DPDNet, that exploits
dual-pixel (DP) views; later work further leveraged DP
data~\cite{abuolaim2021rdpd,pan2021dddnet,xin2021dpdefocus}.
In the more widely applicable single-image setting, KPAC~\cite{son2021kpac}
designs kernel-sharing parallel atrous convolutions tailored to the shape
of defocus blur, IFAN~\cite{lee2021ifan} predicts spatially varying inverse
kernels applied in feature space, DRBNet~\cite{ruan2022drbnet} trains on
light-field-generated and real defocused images with a dynamic residual
design, GKMNet and NRKNet~\cite{quan2021gkmnet,quan2023nrknet} model defocus
with Gaussian kernel mixtures and recursive kernels, and
JDRL~\cite{li2023jdrl} addresses the misalignment inherent to
aperture-bracketed training pairs.
General-purpose restoration backbones such as
Uformer~\cite{wang2022uformer}, Restormer~\cite{zamir2022restormer}, and
NAFNet~\cite{chen2022nafnet}, as well as large-kernel designs such as
LaKDNet~\cite{ruan2023lakdnet}, set strong results on DPDD.
All of these methods are trained on small crops and evaluated at
$\approx2$\,MP, and are therefore implicitly calibrated to the CoC
distribution at that resolution. Our own Stage~1 is a network of this
kind, whose novelty lies not in its low-resolution accuracy but in exposing
a deblur-level map and whole-image affinities for the native-resolution
stage; all other methods can serve as interchangeable anchors through
surrogate signals. Our focus is the step these methods leave open:
producing the native-resolution result.

\paragraph{Training data and defocus synthesis.}
Real training pairs for defocus deblurring are obtained by aperture
bracketing~\cite{abuolaim2020dpdd,lee2021ifan}, which limits datasets to a
few hundred static scenes and introduces misalignment between
captures~\cite{li2023jdrl}. Synthetic alternatives render defocus from
all-in-focus images and depth, using light fields~\cite{ruan2022drbnet},
layered or scattering-based CoC rendering~\cite{peng2022bokehme,
sheng2024dr}, or, most recently, compound-lens models that reproduce
realistic aberrations and bokeh~\cite{compoundlens2026}.
Synthetic data scale easily to high resolution but leave a domain gap to
real optics. We use synthesis to extend scene diversity at native
resolution and fine-tune on real pairs. Orthogonally, we simulate not only
the defocused input but also the \emph{low-resolution anchor}, so that the
upsampler learns to handle the error modes of arbitrary deblurring methods
rather than those of a fixed one; this mirrors degradation modeling in
blind super-resolution~\cite{wang2021realesrgan,zhang2021bsrgan}, applied
here to the output of a restoration method rather than to the sensor.

\paragraph{Generative image restoration and deblurring.}
Diffusion models have become a powerful prior for image restoration.
Early work applied pixel-space diffusion to deblurring and
super-resolution~\cite{saharia2022sr3,whang2022deblurring}; subsequent
methods operate in compact latent spaces~\cite{xia2023diffir,
chen2023hidiff} or adapt large pretrained text-to-image
models~\cite{rombach2022ldm} through control
branches~\cite{wang2024stablesr,lin2024diffbir,yu2024supir,
yang2024pasd,wu2024seesr}.
For defocus specifically, diffusion-based deblurring has been explored with
residual diffusion~\cite{feng2025rddm}, multimodal and depth-aware
guidance~\cite{mmgdiff2025}, blur-level-aware expert routing inside Stable
Diffusion~\cite{mbsd2026}, score-based models for real-world
blur~\cite{scoredefocus2025}, and as the all-in-focus stage of generative
refocusing pipelines~\cite{genrefocus2025}.
The cost of iterative sampling has motivated step
distillation~\cite{salimans2022progressive,song2023consistency,
sauer2024add,yin2024dmd}, and its application to restoration has produced
one- and few-step restorers~\cite{wang2024sinsr,wu2024osediff,xie2024addsr,
onestepdeblur2025}.
These methods demonstrate the value of generative priors for recovering
texture destroyed by blur, but they are designed and evaluated at low or
moderate resolution, and at UHD resolutions they inherit the patch-wise
inference problem discussed below. We follow the same train-then-distill
recipe, but condition the model on the native-resolution observation and a
global low-resolution anchor, and restrict synthesis to regions where the
observation has lost detail.

\paragraph{Ultra-high-definition restoration.}
A growing body of work targets UHD (4K and beyond) restoration, largely for
motion blur, low-light enhancement, and adverse weather; see the recent
survey~\cite{uhdsurvey2025}.
Typical designs process a heavily downsampled version of the input and
recover the full resolution with learned upsampling, bilateral grids, or
frequency-domain operators~\cite{zheng2021uhdmultiguided,
zheng2022cubicmixer,wang2023llformer,wang2024uhdformer}, or adopt
global--local
architectures~\cite{glsgn2022}.
MC-Blur~\cite{zhang2022mcblur} provides a UHD motion-deblurring benchmark,
and most recently an autoregressive flow-matching model performs UHD
deblurring coarse-to-fine by predicting a residual at each
scale~\cite{uhdarflow2026}.
Our decomposition shares the coarse-to-fine spirit of these methods, but
differs in three respects that are specific to defocus: the scale split is
motivated by the growth of the CoC with resolution; the native-resolution
stage explicitly exploits the spatially varying amount of detail that
survives in the defocused observation, rather than predicting the residual
from the coarse estimate alone; and the low-resolution stage is an
arbitrary existing method rather than a jointly trained component.

\paragraph{Patch-wise inference and cross-patch consistency.}
Memory limits force patch-wise processing at UHD resolutions, and it is
well known that restoration networks trained on crops behave differently
when applied to full images; TLC~\cite{chu2022tlc} traces part of this gap
to global statistics aggregation and corrects it at test time.
For generative models, independent patch predictions additionally make
independent stochastic decisions, which produces seams and inconsistent
content.
In high-resolution and panoramic generation, MultiDiffusion and
Mixture-of-Diffusers~\cite{bartal2023multidiffusion,
jimenez2023mixtureofdiffusers} fuse overlapping predictions at every
denoising step, SyncDiffusion~\cite{lee2023syncdiffusion} additionally
synchronizes patches through a perceptual similarity loss,
DemoFusion~\cite{du2024demofusion} and related
works~\cite{he2024scalecrafter} progressively upscale with skip-residual
guidance from a low-resolution generation, and Patch-DM~\cite{ding2023patchdm}
uses feature collage for boundary-free synthesis.
Tiled Diffusion~\cite{madar2025tiled} enforces seamless tiling constraints,
and FrescoDiffusion~\cite{frescodiffusion2026} regularizes tiled denoising
with an upsampled low-resolution latent prior.
In restoration, StableSR~\cite{wang2024stablesr} and SUPIR~\cite{yu2024supir}
aggregate overlapping latent tiles, Patch-Adapter~\cite{patchadapter2025}
combines a low-resolution global adapter with reference-patch attention for
4K inpainting, and OP4KSR~\cite{op4ksr2026} suppresses the periodic
artifacts of one-step patch-free 4K super-resolution.
Two limitations of these mechanisms matter in our setting.
First, most of them were designed for multi-step samplers, where
consistency can be negotiated over many denoising steps; after distillation
to a single step, it must instead be established within one forward pass.
Second, they enforce consistency \emph{locally}, between spatially
overlapping or adjacent tiles, whereas the inconsistencies most visible in
UHD restoration often arise between distant regions that depict the same
content (\eg, two ends of a facade or a textured floor). \method{} addresses
both: its cross-patch attention operates within a single forward pass, and
it is routed by whole-image self-attention computed at low resolution, so
that each patch exchanges information with the related patches wherever
they lie in the image.

\paragraph{Non-local context and spatially adaptive restoration.}
Exploiting recurrence of similar content across an image is a long-standing
principle in restoration, from non-local means and
BM3D~\cite{buades2005nlm,dabov2007bm3d} to learned non-local
operators~\cite{wang2018nonlocal,liu2018nlrn} and sparse or routed
attention that restricts each query to its most relevant
regions~\cite{zhu2023biformer}. More recently, features of self-supervised
and generative foundation models~\cite{oquab2024dinov2,kirillov2023sam}
have been shown to provide dense semantic correspondences without
task-specific training~\cite{amir2021deepvit,tang2023dift}, making them a
robust, training-free source of cross-region affinity.
Computing such affinities directly at tens of megapixels is prohibitive; we
instead obtain them at low resolution from the whole image, combining
learned attention with frozen foundation features, and use them to select
the native-resolution context of each patch, which keeps the per-patch cost
bounded.
Separately, restoration difficulty varies across an image, which has been
exploited to allocate computation adaptively, \eg, by routing patches of
different difficulty to networks of different capacity in
ClassSR~\cite{kong2021classsr}. In defocus deblurring the variation is
especially pronounced and physically grounded, as it follows the defocus
level. \method{} uses a deblur-level map to decide how much to rely on the
observation \vs the generative prior, avoiding both hallucination in sharp
regions and under-restoration in strongly defocused ones.
When this map is not available from the deblurring method, we derive it
from a re-blur agreement surrogate. Re-blurring is a classical tool for
defocus estimation: Zhuo and Sim~\cite{zhuo2011defocusmap} re-blur the
input with a known Gaussian and infer the unknown blur from the ratio of
gradients at edges, and related analysis-by-synthesis
schemes~\cite{karaali2018edgedefocus,lee2019dmenet} select the kernel scale
that best explains the observation.
We apply the same principle to a different pair of images---the
\emph{restored anchor} and the observation---so that the selected kernel
scale measures the deblurring that was performed, and the residual measures
whether the anchor is physically consistent with the input.
Optionally, the disagreement of a method under test-time scale
perturbations serves as an additional uncertainty cue, following standard
practice in ensembles and test-time
augmentation~\cite{lakshminarayanan2017ensembles,ayhan2018tta}, and
monocular depth~\cite{yang2024depthanything,ke2024marigold} provides a
physics-only alternative.

\paragraph{Guided and reference-based upsampling.}
Our upsampler is conditioned on a native-resolution guide, which relates it
to joint-image upsampling. Joint bilateral
upsampling~\cite{kopf2007jbu}, the guided filter~\cite{he2010guided}, and
their learned counterparts~\cite{wu2018deepguided,gharbi2017hdrnet} transfer
edges from a high-resolution guide to a low-resolution solution, and are
widely used for depth and flow upsampling.
Reference-based super-resolution~\cite{zhang2019srntt,yang2020ttsr,
jiang2021c2matching} transfers textures from a separate high-resolution
reference image through correspondence matching.
Both families assume a guide that is sharp, either everywhere or where it is
matched. In our setting the guide is the degraded observation itself, and
its reliability varies continuously across the image with the local amount
of defocus. \method{} models this reliability explicitly, copying
high-frequency content from the guide where it is in focus, attenuating it
as blur increases, and falling back to generative synthesis where it carries
no usable detail.

% =====================================================================
%  Citation keys used (all must be added to the .bib; verify each entry)
% ---------------------------------------------------------------------
%  abuolaim2020dpdd      Abuolaim & Brown, ECCV 2020 (DPDD, DPDNet)
%  abuolaim2021rdpd      Abuolaim et al., ICCV 2021 (multi-task DP / RDPD)
%  pan2021dddnet         Pan et al., CVPR 2021 (dual-pixel deblur + depth)
%  xin2021dpdefocus      Xin et al., ICCV 2021 (defocus map + deblur from DP)
%  son2021kpac           Son et al., ICCV 2021 (KPAC)
%  lee2021ifan           Lee et al., CVPR 2021 (IFAN)
%  ruan2022drbnet        Ruan et al., CVPR 2022 (DRBNet)
%  ruan2023lakdnet       Ruan et al., 2023 (LaKDNet)
%  quan2021gkmnet        Quan et al., ICCV 2021 (Gaussian kernel mixture)
%  quan2023nrknet        Quan et al., CVPR 2023 (NRKNet)
%  li2023jdrl            Li et al., AAAI 2023 (misaligned training pairs)
%  zamir2022restormer    Zamir et al., CVPR 2022
%  wang2022uformer       Wang et al., CVPR 2022
%  chen2022nafnet        Chen et al., ECCV 2022
%  zhuo2011defocusmap    Zhuo & Sim, PR 2011
%  karaali2018edgedefocus Karaali & Jung, TIP 2018
%  lee2019dmenet         Lee et al., CVPR 2019 (DMENet)
%  krishnan2009hyperlaplacian Krishnan & Fergus, NeurIPS 2009
%  saharia2022sr3        Saharia et al., TPAMI 2022 (SR3)
%  whang2022deblurring   Whang et al., CVPR 2022
%  xia2023diffir         Xia et al., ICCV 2023
%  chen2023hidiff        Chen et al., NeurIPS 2023 (HI-Diff)
%  rombach2022ldm        Rombach et al., CVPR 2022
%  wang2024stablesr      Wang et al., IJCV 2024
%  lin2024diffbir        Lin et al., ECCV 2024
%  yu2024supir           Yu et al., CVPR 2024
%  yang2024pasd          Yang et al., ECCV 2024
%  wu2024seesr           Wu et al., CVPR 2024
%  wang2024sinsr         Wang et al., CVPR 2024
%  wu2024osediff         Wu et al., NeurIPS 2024
%  xie2024addsr          Xie et al., 2024 (AddSR)
%  onestepdeblur2025     One-Step Diffusion for Motion Deblurring, arXiv 2503.06537
%  feng2025rddm          Feng et al., AAAI 2025
%  mmgdiff2025           Multimodal-guided diffusion + depth-aware fusion, PR 2025
%  mbsd2026              Multi-Level Blur-Aware SD, AAAI 2026
%  scoredefocus2025      Score-based real-world defocus, Sci. Rep. 2025
%  genrefocus2025        Generative Refocusing, arXiv 2512.16923
%  liang2021swinir       Liang et al., ICCVW 2021
%  wang2021realesrgan    Wang et al., ICCVW 2021
%  uhdsurvey2025         UHD restoration survey, arXiv 2505.16161
%  zheng2021uhdmultiguided Zheng et al., CVPR 2021 (UHD multi-guided bilateral)
%  zheng2022cubicmixer   Zheng et al., arXiv 2206.03678
%  wang2023llformer      Wang et al., AAAI 2023 (UHD-LOL / LLFormer)
%  wang2024uhdformer     Wang et al., AAAI 2024 (UHDformer)
%  glsgn2022             Global-Local Stepwise GAN, arXiv 2207.08808
%  zhang2022mcblur       Zhang et al., MC-Blur, TCSVT 2023 / arXiv 2112.00234
%  uhdarflow2026         UHD deblurring via AR flow, arXiv 2603.10517
%  chu2022tlc            Chu et al., ECCV 2022 (TLC)
%  bartal2023multidiffusion Bar-Tal et al., ICML 2023
%  jimenez2023mixtureofdiffusers Jimenez, arXiv 2023
%  lee2023syncdiffusion  Lee et al., NeurIPS 2023
%  du2024demofusion      Du et al., CVPR 2024
%  he2024scalecrafter    He et al., ICLR 2024
%  ding2023patchdm       Ding et al., WACV 2023 (Patch-DM)
%  madar2025tiled        Madar & Fried, CVPR 2025
%  frescodiffusion2026   FrescoDiffusion, arXiv 2603.17555
%  patchadapter2025      Patch-Adapter, arXiv 2510.13419
%  op4ksr2026            OP4KSR, arXiv 2605.13457
%  kopf2007jbu           Kopf et al., SIGGRAPH 2007
%  he2010guided          He et al., ECCV 2010 / TPAMI 2013
%  wu2018deepguided      Wu et al., CVPR 2018
%  gharbi2017hdrnet      Gharbi et al., SIGGRAPH 2017
%  zhang2019srntt        Zhang et al., CVPR 2019
%  yang2020ttsr          Yang et al., CVPR 2020
%  jiang2021c2matching   Jiang et al., CVPR 2021
%  salimans2022progressive Salimans & Ho, ICLR 2022 (progressive distillation)
%  song2023consistency   Song et al., ICML 2023 (consistency models)
%  sauer2024add          Sauer et al., ECCV 2024 (adversarial diffusion distillation)
%  yin2024dmd            Yin et al., CVPR 2024 (distribution matching distillation)
%  buades2005nlm         Buades et al., CVPR 2005 (non-local means)
%  dabov2007bm3d         Dabov et al., TIP 2007 (BM3D)
%  wang2018nonlocal      Wang et al., CVPR 2018 (non-local neural networks)
%  liu2018nlrn           Liu et al., NeurIPS 2018 (NLRN)
%  zhu2023biformer       Zhu et al., CVPR 2023 (BiFormer, bi-level routing attention)
%  kong2021classsr       Kong et al., CVPR 2021 (ClassSR)
%  oquab2024dinov2       Oquab et al., TMLR 2024 (DINOv2)
%  kirillov2023sam       Kirillov et al., ICCV 2023 (Segment Anything)
%  amir2021deepvit       Amir et al., ECCVW 2022 / arXiv 2021 (Deep ViT features as dense descriptors)
%  tang2023dift          Tang et al., NeurIPS 2023 (DIFT, diffusion features for correspondence)
%  peng2022bokehme       Peng et al., CVPR 2022 (BokehMe)
%  sheng2024dr           Sheng et al., CVPR 2024 (Dr.Bokeh)  -- VERIFY venue/key
%  compoundlens2026      Realistic Compound-Lens Defocus Blur Synthesis, arXiv 2607.05837
%  zhang2021bsrgan       Zhang et al., ICCV 2021 (BSRGAN degradation model)
%  yang2024depthanything Yang et al., CVPR 2024 (Depth Anything)
%  ke2024marigold        Ke et al., CVPR 2024 (Marigold)
%  lakshminarayanan2017ensembles Lakshminarayanan et al., NeurIPS 2017 (deep ensembles)
%  ayhan2018tta          Ayhan & Berens, MIDL 2018 (test-time augmentation uncertainty)
% =====================================================================

% =====================================================================
%  Method subsection: surrogate interface signals for off-the-shelf
%  low-resolution deblurring methods ("cross play").
%  Uses macros from intro_related.tex (\method, \todo, \eg, \ie).
%  Requires: amsmath, amssymb
% =====================================================================

\subsection{Surrogate Interface for Off-the-Shelf Methods}
\label{sec:surrogates}

Stage~2 consumes three signals from Stage~1: the anchor $\hat{x}$, the
deblur-level map $m$, and the affinity map $A$
(Sec.~\todo{ref to interface section}).
Our own Stage~1 produces $m$ and $A$ natively. An arbitrary deblurring
method $\mathcal{M}$, in contrast, is a black box that returns only
$\hat{x} = \mathcal{M}(y)$, where $y$ denotes the defocused input
downsampled by the factor $s=4$. We construct $m$ and $A$ from quantities
that can be measured for any such method.

\paragraph{Notation.}
All quantities in this section are computed at low resolution.
$p$ indexes pixels, $G_\sigma$ is a Gaussian window of standard deviation
$\sigma$ that defines local neighborhoods, and $\lvert\cdot\rvert$ is
applied per pixel and averaged over color channels.
$k_r$ denotes a normalized defocus kernel of radius $r$, and
$\mathcal{R} = \{r_0{=}0, r_1, \dots, r_{K}\}$ is a discrete set of radii
spanning the CoC range observed at low resolution (\todo{values}).

\paragraph{Re-blur agreement.}
If $\hat{x}$ is a faithful all-in-focus estimate, then re-blurring it with
the correct local defocus kernel should reproduce the observation, since
defocus is well approximated locally by a convolution with a kernel whose
size depends on depth~\cite{zhuo2011defocusmap}. For each radius
$r \in \mathcal{R}$ we therefore compute a local, contrast-normalized
re-blur error
\begin{equation}
  e_r(p) \;=\;
  \frac{\big(G_\sigma * \lvert k_r * \hat{x} - y \rvert\big)(p)}
       {c(p) + \epsilon},
  \qquad
  c(p) = \sqrt{\big(G_\sigma * (y - G_\sigma * y)^2\big)(p)},
  \label{eq:reblur_error}
\end{equation}
where $c$ is the local standard deviation of $y$, which makes errors
comparable between flat and textured regions and reduces the influence of
noise. From the error profile over $\mathcal{R}$ we extract
\begin{align}
  r^*(p) &= \arg\min_{r \in \mathcal{R}} e_r(p),
  \label{eq:rstar}\\
  \rho(p) &= e_{r^*(p)}(p).
  \label{eq:rho}
\end{align}
The selected radius $r^*$ estimates \emph{how much deblurring was
performed} at $p$: it is close to zero in regions that were already in
focus and grows with the defocus that $\mathcal{M}$ removed. The residual
$\rho$ measures \emph{reliability}: when no kernel in $\mathcal{R}$ explains
$y$ from $\hat{x}$, the anchor is inconsistent with the observation, \eg,
because of hallucinated structure, ringing, or tone shifts.
In practice we replace the hard $\arg\min$ by a soft-argmin over
$\mathcal{R}$ with temperature $\tau$ for sub-radius precision, and use a
disc--Gaussian mixture for $k_r$ to account for the non-ideal shape of real
defocus kernels (\todo{fit on DPDD train?}).
The sweep requires $K{+}1$ convolutions at low resolution, which is
negligible compared with Stage~2.

\paragraph{Unrestored regions.}
Re-blur agreement alone has a blind spot. If $\mathcal{M}$ fails to
restore a strongly defocused region and returns it essentially unchanged,
$\hat{x} \approx y$ there, so that $r^* \approx 0$ and $\rho \approx 0$:
the region appears both in focus and reliable. We resolve this ambiguity
with a no-reference sharpness measure of the input,
\begin{equation}
  s(p) = \frac{\big(G_\sigma * \lvert \nabla y \rvert\big)(p)}
              {c(p) + \epsilon},
  \label{eq:sharpness}
\end{equation}
which is low wherever the observation lacks high-frequency content,
regardless of what $\mathcal{M}$ did. A region with small $r^*$ \emph{and}
small $s$ is thus identified as defocused but unrestored, and hence as
requiring generation.

\paragraph{Optional: perturbation disagreement.}
Where additional low-resolution passes are affordable, we measure the
instability of $\mathcal{M}$ under input rescaling,
\begin{equation}
  v(p) = \operatorname{Var}_{\alpha \in \mathcal{A}}
         \Big[\, \mathcal{U}_{1/\alpha}\big(\mathcal{M}(\mathcal{U}_{\alpha}\, y)\big)(p) \Big],
  \label{eq:tta}
\end{equation}
where $\mathcal{U}_\alpha$ resamples by factor $\alpha$ and
$\mathcal{A} = \{0.9, 1, 1.1\}$ (\todo{values}). We use scale rather than
flip perturbations, to which many restoration networks are nearly
equivariant. Disagreement is a standard proxy for predictive
uncertainty~\cite{lakshminarayanan2017ensembles,ayhan2018tta} and
complements $\rho$, which detects inconsistency with the observation but
not ambiguity among plausible restorations.

\paragraph{Calibration to the native map.}
Stage~2 is trained on the native map $m$ produced by our Stage~1. Rather
than fusing the surrogate cues by hand, we learn a lightweight calibration
head $h_\phi$, a small convolutional network with a local receptive field
(\todo{architecture}), that maps them to the same format:
\begin{equation}
  \tilde{m} = h_\phi\big(\, r^*/r_K,\; \rho,\; s,\; [v] \,\big),
  \qquad
  \phi^\star = \arg\min_\phi\;
  \mathbb{E}_{\mathcal{M} \sim \mathcal{P},\, y}
  \big\lVert h_\phi(\cdot) - m \big\rVert_1 ,
  \label{eq:calib}
\end{equation}
where $[v]$ is included only when perturbation disagreement is used, $m$ is
the native map of our Stage~1 for the same input, and the expectation is
over training images and a pool $\mathcal{P}$ of deblurring methods
together with our simulated anchors. Because every input to $h_\phi$ is a
black-box measurement, the calibrated surrogate remains method-agnostic,
and we evaluate it on methods excluded from $\mathcal{P}$.
The surrogate map $\tilde{m}$ is bilinearly upsampled to native resolution
and cropped per patch in the same way as $m$.

\paragraph{Affinity surrogate.}
For the affinity map, we extract features $F = \Phi(\hat{x})$ from a frozen
vision foundation model $\Phi$~\cite{oquab2024dinov2} on the anchor, pool
them to the patch grid used by Stage~2, and define
\begin{equation}
  \tilde{A}_{ij} = \operatorname{softmax}_j
  \Big( \tfrac{1}{\sqrt{d}}\, \bar{F}_i^{\top} \bar{F}_j \Big),
  \label{eq:affinity}
\end{equation}
where $\bar{F}_i \in \mathbb{R}^d$ is the $\ell_2$-normalized pooled feature
of patch $i$. As with the native affinity map $A$, the related patches
for patch $i$ are the top-$k$ entries of row $i$ (\todo{$k$}).

\paragraph{Training with interchangeable signals.}
To make Stage~2 robust to which variant of each signal it receives, we
replace, during training, $m$ by $\tilde{m}$ and $A$ by $\tilde{A}$
independently with probability $p_{\text{sur}}$ (\todo{value}).
In Sec.~\todo{ref experiments} we report the gap between native and
surrogate signals for our own Stage~1, and the gains obtained with
surrogate signals for off-the-shelf methods.

\paragraph{Limitations of the surrogate.}
The local convolution model underlying Eq.~\eqref{eq:reblur_error} is
violated at occlusion boundaries, where foreground and background blur mix;
we keep $\sigma$ small to confine the resulting errors. Sensor noise
inflates $e_r$ uniformly across radii and is largely absorbed by the
contrast normalization. Finally, the surrogate cannot distinguish a
correct restoration from an incorrect one that is equally consistent with
the observation, since both re-blur to $y$; the perturbation disagreement
$v$ partially addresses this case.